\chapter{Our Toolbox}
\label{ch:toolbox}

This chapter describes the two tools we rely on: the Soot framework for static analysis of Java
bytecode, and OpenJ9's debug counters for measuring what the JIT compiler does.

\section{Soot}

Soot~\cite{vallee2010soot} is a framework for analysing and transforming Java bytecode. It lifts
bytecode into \emph{Jimple}, a typed three-address intermediate representation, performs
analyses and optimizations on Jimple, and writes the optimized program back as bytecode. The
static escape analysis used by optimistic stack allocation~\cite{anand2024optimistic} and the
static half of PYE~\cite{thakur2019pye} are both implemented in Soot.

\paragraph{Data-flow analysis framework.} Soot provides a generic data-flow analysis framework:
\begin{itemize}
  \item it is template driven: the analysis writer supplies the flow set, the flow functions, the
        merge function and the direction (forward or backward);
  \item Soot then computes the fixed point itself, using worklists; and
  \item it provides further utilities such as def-use chains, points-to analysis and call-graph
        construction.
\end{itemize}

\subsection{Why a Three-Address Code?}

Java bytecode is a stack-based instruction set: most instructions take their operands from an
implicit operand stack and push their results back onto it. This hides the data flow.
Table~\ref{tab:tac} shows the bytecode of a small method that creates a \texttt{Point} and
returns the sum of its coordinates, the state of the operand stack after each instruction, and
the corresponding three-address code (TAC).

In bytecode the operands are implicit, so to find out what \texttt{iadd} adds, or which
\texttt{new} reaches the local \texttt{p}, an analysis must simulate the operand stack. In TAC
every statement performs one operation on explicit operands, so the data flow can be read
directly. This is why Soot's analyses, including escape analysis, work on Jimple rather than on
bytecode.

\section{OpenJ9 Debug Counters}

To understand what the JIT compiler does on real programs, we use OpenJ9's \emph{debug counters},
a utility for counting arbitrary events inside the JIT compiler. There are two kinds:
\begin{itemize}
  \item \textbf{Static} counters are incremented at compile time; nothing is added to the
        generated code. They count compile-time events, such as how many allocation sites were
        rejected for a given reason.
  \item \textbf{Dynamic} counters insert counter increments into the generated code, so they
        count how often an event happens at run time.
\end{itemize}

Counter names are hierarchical: a `\texttt{/}' in a name works like a slash in a file path and
expresses containment. When the counters are printed, each parent counter shows the total of its
children, and each child also shows its percentage of that parent. Listing~\ref{lst:dbg} and
Table~\ref{tab:dbg} show an example.

\begin{lstlisting}[language=C++, caption={Counting inlined and non-inlined call sites.}, label={lst:dbg}]
prependDebugCounter("callSites/inlined", tt);
...
prependDebugCounter("callSites/notInlined", tt);
\end{lstlisting}

\begin{table}[htbp]
  \centering
  \begin{tabular}{@{}rlrr@{}}
    \toprule
    & Counter & Count & \% \\
    \midrule
    1: & \texttt{callSites}             & 10000 & -- \\
    2: & \texttt{callSites/inlined}     & 8000  & 80.00\% \\
    3: & \texttt{callSites/notInlined}  & 2000  & 20.00\% \\
    \bottomrule
  \end{tabular}
  \caption{Output produced by the counters in Listing~\ref{lst:dbg}.}
  \label{tab:dbg}
\end{table}

Both experiments in Chapter~\ref{ch:experiments} use static debug counters.
